% ═══════════════════════════════════════════════════════════════════════════════
%  ANNEXES
% ═══════════════════════════════════════════════════════════════════════════════
\newpage
\appendix
\section*{Annexes}
\addcontentsline{toc}{section}{Annexes}

% ── Annexe A : Glossaire ──────────────────────────────────────────────────────
\section{Glossaire technique}
\label{annex:glossaire}

\begin{xltabular}{\textwidth}{lX}
    \toprule
    \textbf{Terme} & \textbf{Définition} \\
    \midrule
    \endfirsthead

    \toprule
    \textbf{Terme} & \textbf{Définition} \\
    \midrule
    \endhead

    \bottomrule
    \endfoot

    JWT & \textit{JSON Web Token} -- standard ouvert (RFC 7519) pour l'échange sécurisé d'informations entre parties sous forme d'objet JSON signé numériquement. Dans ce système, les JWT sont stockés dans un cookie \texttt{httpOnly} pour l'authentification stateless. \\
    RBAC & \textit{Role-Based Access Control} -- modèle de contrôle d'accès où les permissions sont assignées à des rôles plutôt qu'à des utilisateurs individuels. Le système utilise quatre rôles : administrateur, coordinateur, enseignant, étudiant. \\
    SPA & \textit{Single Page Application} -- architecture frontend où une seule page HTML est chargée initialement, puis le contenu est mis à jour dynamiquement via JavaScript sans rechargement complet de la page. \\
    AOP & \textit{Aspect-Oriented Programming} -- paradigme de programmation permettant d'isoler des préoccupations transversales (logging, audit, sécurité) dans des modules autonomes appelés aspects. Utilisé dans ce système pour le journal d'audit. \\
    DTO & \textit{Data Transfer Object} -- objet de transfert de données servant à encapsuler les données échangées entre les couches de l'architecture (contrôleur, service, repository) sans exposer les entités JPA directement. \\
    MLD & \textit{Modèle Logique de Données} -- représentation structurée des données d'une base relationnelle, dérivée du modèle conceptuel (MCD) et traduisant les entités en tables avec leurs clés et relations. \\
    CI/CD & \textit{Continuous Integration / Continuous Deployment} -- ensemble de pratiques automatisant la construction, les tests et le déploiement du code à chaque modification. Implémenté via GitHub Actions. \\
    DnD & \textit{Drag and Drop} (glisser-déposer) -- interaction utilisateur permettant de déplacer des éléments visuels (jurys) vers des emplacements cibles (créneaux du planning) dans le \textit{Defense Designer}. \\
    ORM & \textit{Object-Relational Mapping} -- technique de mapping entre les objets d'un langage orienté objet (Java) et les tables d'une base de données relationnelle. Implémenté via Hibernate/JPA. \\
    API REST & \textit{Application Programming Interface -- Representational State Transfer} -- style d'architecture pour les interfaces web utilisant les verbes HTTP (GET, POST, PUT, PATCH, DELETE) pour les opérations CRUD. \\
    CORS & \textit{Cross-Origin Resource Sharing} -- mécanisme de sécurité permettant de contrôler les requêtes inter-domaines. Configuré pour autoriser les origines du frontend en développement et production. \\
    XSS & \textit{Cross-Site Scripting} -- type de vulnerability permettant d'injecter du code JavaScript dans les pages consultées par d'autres utilisateurs. Prévenu par les cookies \texttt{httpOnly} et la validation des entrées. \\
    BCrypt & Algorithme de hachage de mots de passe conçu pour être résistant aux attaques par force brute. Utilisé pour stocker les mots de passe de manière sécurisée. \\
    OpenHTMLToPDF & Bibliothèque Java open source convertissant du HTML et CSS en PDF. Basée sur Flying Saucer, elle est utilisée pour la génération côté serveur des documents administratifs. \\
    Thymeleaf & Moteur de templates Java permettant de générer du HTML côté serveur. Utilisé pour construire les modèles de documents PDF avant conversion via OpenHTMLToPDF. \\
    WebSocket & Protocole de communication bidirectionnel en temps réel entre le client et le serveur. Mentionné dans les perspectives d'évolution pour les notifications temps réel. \\
    PWA & \textit{Progressive Web Application} -- application web utilisant des technologies modernes pour offrir une expérience similaire à une application native, incluant le mode hors ligne et les notifications push. \\
    LMD & \textit{Licence -- Master -- Doctorat} -- système d'organisation des études supérieur adopté au Maroc, structurant les formations en trois cycles successifs. \\
    PlantUML & Outil de modélisation UML basé sur du texte, permettant de générer automatiquement des diagrammes à partir de descriptions semi-formelles. \\

    \bottomrule
\end{xltabular}

% ── Annexe B : Endpoints API ─────────────────────────────────────────────────
\section{Liste complète des endpoints \acrshort{api}}
\label{annex:api_endpoints}

L'\acrshort{api} \acrshort{rest} expose 124 endpoints organisés par module. Le préfixe commun est \texttt{/api}.

\subsection{Authentification}

\begin{xltabular}{\textwidth}{llX}
    \toprule
    \textbf{Méthode} & \textbf{Chemin} & \textbf{Description} \\
    \midrule
    POST  & /auth/login              & Authentifier un utilisateur (cookie JWT) \\
    POST  & /auth/logout             & Déconnexion (effacement du cookie) \\
    POST  & /auth/forgot-password    & Demander un lien de réinitialisation \\
    POST  & /auth/reset-password     & Réinitialiser le mot de passe via un token \\
    POST  & /auth/verify-account     & Vérifier le compte via le token d'activation \\
    \bottomrule
\end{xltabular}

\subsection{Profil utilisateur}

\begin{xltabular}{\textwidth}{llX}
    \toprule
    \textbf{Méthode} & \textbf{Chemin} & \textbf{Description} \\
    \midrule
    GET    & /me                     & Obtenir le profil de l'utilisateur courant \\
    PATCH  & /me                     & Modifier le prénom et/ou le nom \\
    PUT    & /me/password            & Changer le mot de passe \\
    \bottomrule
\end{xltabular}

\subsection{Administration -- Gestion des utilisateurs}

\begin{xltabular}{\textwidth}{llX}
    \toprule
    \textbf{Méthode} & \textbf{Chemin} & \textbf{Description} \\
    \midrule
    GET     & /admin/users                & Lister les utilisateurs (paginé, filtrable par rôle) \\
    POST    & /admin/users                & Créer un utilisateur \\
    POST    & /admin/users/bulk           & Création massive d'utilisateurs \\
    PUT     & /admin/users/\{id\}          & Modifier un utilisateur \\
    DELETE  & /admin/users/\{id\}          & Supprimer un utilisateur \\
    GET     & /coordinator/users          & Lister les utilisateurs (par rôle STUDENT/TEACHER) \\
    \bottomrule
\end{xltabular}

\subsection{Administration -- Configuration académique}

\begin{xltabular}{\textwidth}{llX}
    \toprule
    \textbf{Méthode} & \textbf{Chemin} & \textbf{Description} \\
    \midrule
    GET     & /admin/faculties              & Lister les facultés \\
    POST    & /admin/faculties              & Créer une faculté \\
    GET     & /admin/faculties/\{id\}        & Détails d'une faculté \\
    PUT     & /admin/faculties/\{id\}        & Modifier une faculté \\
    PATCH   & /admin/faculties/\{id\}        & Modification partielle d'une faculté \\
    DELETE  & /admin/faculties/\{id\}        & Supprimer une faculté \\
    \midrule
    GET     & /admin/departments            & Lister les départements \\
    POST    & /admin/departments            & Créer un département \\
    GET     & /admin/departments/\{id\}      & Détails d'un département \\
    PUT     & /admin/departments/\{id\}      & Modifier un département \\
    PATCH   & /admin/departments/\{id\}      & Modification partielle d'un département \\
    DELETE  & /admin/departments/\{id\}      & Supprimer un département \\
    \midrule
    GET     & /admin/config/majors          & Lister les filières \\
    POST    & /admin/config/majors          & Créer une filière \\
    PUT     & /admin/config/majors/\{id\}    & Modifier une filière \\
    PATCH   & /admin/config/majors/\{id\}    & Modification partielle d'une filière \\
    DELETE  & /admin/config/majors/\{id\}    & Supprimer une filière \\
    \midrule
    GET     & /admin/config/levels          & Lister les niveaux \\
    POST    & /admin/config/levels          & Créer un niveau \\
    PUT     & /admin/config/levels/\{id\}    & Modifier un niveau \\
    PATCH   & /admin/config/levels/\{id\}    & Modification partielle d'un niveau \\
    DELETE  & /admin/config/levels/\{id\}    & Supprimer un niveau \\
    \midrule
    GET     & /admin/config/teacher-ranks     & Lister les grades \\
    POST    & /admin/config/teacher-ranks     & Créer un grade \\
    PUT     & /admin/config/teacher-ranks/\{id\} & Modifier un grade \\
    PATCH   & /admin/config/teacher-ranks/\{id\} & Modification partielle d'un grade \\
    DELETE  & /admin/config/teacher-ranks/\{id\} & Supprimer un grade \\
    \midrule
    GET     & /admin/config/jury-role-templates     & Lister les templates de rôles de jury \\
    POST    & /admin/config/jury-role-templates     & Créer un template \\
    PUT     & /admin/config/jury-role-templates/\{id\} & Modifier un template \\
    DELETE  & /admin/config/jury-role-templates/\{id\} & Supprimer un template \\
    \midrule
    GET     & /admin/config/settings        & Obtenir les paramètres de soutenance \\
    PUT     & /admin/config/settings        & Modifier les paramètres (complet) \\
    PATCH   & /admin/config/settings        & Modifier les paramètres (partiel) \\
    \bottomrule
\end{xltabular}

\subsection{Administration -- Salles, Statistiques, Audit}

\begin{xltabular}{\textwidth}{llX}
    \toprule
    \textbf{Méthode} & \textbf{Chemin} & \textbf{Description} \\
    \midrule
    GET     & /admin/rooms                  & Lister les salles \\
    POST    & /admin/rooms                  & Créer une salle \\
    POST    & /admin/rooms/bulk             & Création massive de salles \\
    PUT     & /admin/rooms/\{id\}            & Modifier une salle \\
    PATCH    & /admin/rooms/\{id\}            & Modification partielle d'une salle \\
    DELETE  & /admin/rooms/\{id\}            & Supprimer une salle \\
    \midrule
    GET     & /admin/stats                  & Statistiques globales \\
    \midrule
    GET     & /admin/audit-logs             & Journal d'audit (paginé) \\
    POST    & /admin/audit-logs             & Créer une entrée d'audit manuellement \\
    \bottomrule
\end{xltabular}

\subsection{Coordinateur -- Sessions de soutenance}

\begin{xltabular}{\textwidth}{llX}
    \toprule
    \textbf{Méthode} & \textbf{Chemin} & \textbf{Description} \\
    \midrule
    GET     & /coordinator/defense-sessions                  & Lister les sessions \\
    POST    & /coordinator/defense-sessions                  & Créer une session \\
    PUT     & /coordinator/defense-sessions/\{id\}            & Modifier une session \\
    DELETE  & /coordinator/defense-sessions/\{id\}            & Supprimer une session \\
    POST    & /coordinator/defense-sessions/\{id\}/transition & Changer le statut \\
    PATCH   & /coordinator/defense-sessions/\{id\}/freeze     & Geler la session \\
    PATCH   & /coordinator/defense-sessions/\{id\}/unfreeze   & Dégeler la session \\
    PATCH   & /coordinator/defense-sessions/\{id\}/approve    & Approuver la session (ADMIN) \\
    PATCH   & /coordinator/defense-sessions/\{id\}/revoke-approval & Retirer l'approbation \\
    \bottomrule
\end{xltabular}

\subsection{Coordinateur -- Projets, Groupes, Jurys}

\begin{xltabular}{\textwidth}{llX}
    \toprule
    \textbf{Méthode} & \textbf{Chemin} & \textbf{Description} \\
    \midrule
    GET     & /coordinator/projects              & Lister les projets \\
    POST    & /coordinator/projects              & Créer un projet \\
    POST    & /coordinator/projects/bulk         & Import massive de projets \\
    PUT     & /coordinator/projects/\{id\}        & Modifier un projet \\
    PATCH   & /coordinator/projects/\{id\}/status & Changer le statut du projet \\
    DELETE  & /coordinator/projects/\{id\}        & Supprimer un projet \\
    \midrule
    GET     & /coordinator/groups                & Lister les groupes \\
    POST    & /coordinator/groups                & Créer un groupe \\
    DELETE  & /coordinator/groups/\{id\}          & Supprimer un groupe \\
    DELETE  & /coordinator/groups/\{id\}/members/\{studentId\} & Retirer un membre \\
    \midrule
    GET     & /coordinator/juries                & Lister les jurys \\
    POST    & /coordinator/juries                & Créer un jury \\
    PUT     & /coordinator/juries/\{id\}          & Modifier un jury \\
    DELETE  & /coordinator/juries/\{id\}          & Supprimer un jury \\
    \bottomrule
\end{xltabular}

\subsection{Coordinateur -- Planning et Conflits}

\begin{xltabular}{\textwidth}{llX}
    \toprule
    \textbf{Méthode} & \textbf{Chemin} & \textbf{Description} \\
    \midrule
    GET     & /coordinator/schedules                 & Obtenir le planning courant \\
    POST    & /coordinator/schedules                 & Sauvegarder le planning \\
    POST    & /coordinator/schedules/generation      & Auto-génération du planning \\
    PATCH   & /coordinator/schedules/publication     & Publier le planning \\
    \midrule
    POST    & /coordinator/conflicts/validate        & Valider le planning contre les conflits \\
    \midrule
    GET     & /coordinator/grades                    & Consulter les notes et moyennes \\
    GET     & /coordinator/unavailability            & Consulter les indisponibilités \\
    GET     & /coordinator/stats                     & Statistiques du coordinateur \\
    \bottomrule
\end{xltabular}

\subsection{Coordinateur -- Défenses}

\begin{xltabular}{\textwidth}{llX}
    \toprule
    \textbf{Méthode} & \textbf{Chemin} & \textbf{Description} \\
    \midrule
    POST    & /coordinator/defenses/\{id\}/cancel    & Annuler une soutenance planifiée \\
    \bottomrule
\end{xltabular}

\subsection{Coordinateur -- Documents}

\begin{xltabular}{\textwidth}{llX}
    \toprule
    \textbf{Méthode} & \textbf{Chemin} & \textbf{Description} \\
    \midrule
    POST  & /coordinator/documents/pdf/evaluation-sheets      & Générer la fiche d'évaluation PDF \\
    POST  & /coordinator/documents/pdf/attendance-lists       & Générer la liste de présence PDF \\
    POST  & /coordinator/documents/pdf/jury-convocations      & Générer la convocation jury PDF \\
    POST  & /coordinator/documents/pdf/schedule               & Générer le planning PDF \\
    POST  & /coordinator/documents/pdf/proces-verbal          & Générer le PV PDF \\
    \midrule
    POST  & /coordinator/documents/evaluation-sheets      & Données de la fiche d'évaluation \\
    POST  & /coordinator/documents/attendance-lists       & Données de la liste de présence \\
    POST  & /coordinator/documents/jury-convocations      & Données de la convocation jury \\
    POST  & /coordinator/documents/schedule               & Données du planning imprimable \\
    POST  & /coordinator/documents/proces-verbal          & Données du procès-verbal \\
    \midrule
    POST  & /coordinator/documents/evaluation-sheets/pdf  & Télécharger la fiche d'évaluation PDF \\
    POST  & /coordinator/documents/attendance-lists/pdf   & Télécharger la liste de présence PDF \\
    POST  & /coordinator/documents/jury-convocations/pdf  & Télécharger la convocation jury PDF \\
    POST  & /coordinator/documents/schedule/pdf           & Télécharger le planning PDF \\
    POST  & /coordinator/documents/proces-verbal/pdf      & Télécharger le PV PDF \\
    \bottomrule
\end{xltabular}

\subsection{Enseignant}

\begin{xltabular}{\textwidth}{llX}
    \toprule
    \textbf{Méthode} & \textbf{Chemin} & \textbf{Description} \\
    \midrule
    GET   & /teacher/evaluations               & Lister les évaluations assignées \\
    POST  & /teacher/evaluations/\{id\}         & Soumettre une évaluation (score + commentaire) \\
    \midrule
    GET   & /teacher/unavailabilities          & Obtenir les créneaux d'indisponibilité \\
    POST  & /teacher/unavailabilities          & Enregistrer les créneaux d'indisponibilité \\
    \midrule
    GET   & /teacher/schedules                 & Obtenir le planning personnel \\
    \midrule
    GET   & /teacher/stats                     & Statistiques personnelles \\
    \bottomrule
\end{xltabular}

\subsection{Étudiant}

\begin{xltabular}{\textwidth}{llX}
    \toprule
    \textbf{Méthode} & \textbf{Chemin} & \textbf{Description} \\
    \midrule
    GET     & /student/groups                 & Obtenir l'espace groupe \\
    POST    & /student/groups                 & Créer un groupe \\
    POST    & /student/groups/\{id\}/members   & Rejoindre un groupe \\
    DELETE  & /student/groups/leave            & Quitter son groupe \\
    \midrule
    GET     & /student/documents              & Lister les documents \\
    POST    & /student/documents/\{id\}/attachments & Téléverser une pièce jointe \\
    GET     & /student/documents/\{id\}/download    & Télécharger un document \\
    \midrule
    GET     & /student/defenses               & Obtenir les infos de sa soutenance \\
    GET     & /student/convocations           & Télécharger sa convocation PDF \\
    \midrule
    GET     & /student/stats                  & Statistiques personnelles \\
    \bottomrule
\end{xltabular}

\subsection{Notifications}

\begin{xltabular}{\textwidth}{llX}
    \toprule
    \textbf{Méthode} & \textbf{Chemin} & \textbf{Description} \\
    \midrule
    GET    & /notifications                  & Lister les notifications (paginé) \\
    PATCH  & /notifications/\{id\}/read       & Marquer une notification comme lue \\
    PATCH  & /notifications/read-all          & Marquer toutes les notifications comme lues \\
    POST   & /notifications/\{id\}/send-email & Envoyer un email associé à une notification \\
    \bottomrule
\end{xltabular}

% ── Annexe C : Manuel d'utilisation rapide ────────────────────────────────────
\section{Manuel d'utilisation rapide}
\label{annex:manuel}

Ce manuel décrit les étapes essentielles pour chaque rôle de l'application.

\subsection{Administrateur}

\textbf{Première connexion et configuration initiale :}
\begin{enumerate}
    \item Se connecter avec les identifiants administrateur fournis.
    \item Créer les départements, les filières, les niveaux et les salles.
    \item Créer les comptes des coordinateurs, enseignants et étudiants (individuellement ou en masse via l'import Excel).
\end{enumerate}

\textbf{Tâches récurrentes :}
\begin{itemize}
    \item Consulter le tableau de bord (\texttt{/admin}) pour visualiser les statistiques globales.
    \item Gérer les comptes utilisateurs (création, modification, désactivation).
    \item Consulter le journal d'audit (\texttt{/admin/audit-logs}) pour tracer les actions effectuées.
    \item Approuver les sessions de soutenance (\texttt{approve}) avant leur publication.
\end{itemize}

\subsection{Coordinateur}

\textbf{Préparation d'une session de soutenance :}
\begin{enumerate}
    \item Créer une session de soutenance (\texttt{/coordinator/defense-sessions}) en définissant les dates, les coefficients d'évaluation et la session active.
    \item Importer ou créer les projets (\texttt{/coordinator/projects}) avec les informations des étudiants et encadrants.
    \item Constituer les jurys (\texttt{/coordinator/juries}) en assignant les enseignants aux rôles (Président, Rapporteur, Examinateur).
    \item Activer la session (\textit{transition} DRAFT $\rightarrow$ ACTIVE).
\end{enumerate}

\textbf{Planification des soutenances :}
\begin{enumerate}
    \item Accéder au \textit{Defense Designer} (\texttt{/coordinator/schedule}).
    \item Glisser-déposer les jurys depuis la barre latérale vers les créneaux du calendrier.
    \item En cas de conflit, le moteur de détection affiche un message explicatif avec suggestion de résolution.
    \item Optionnel : utiliser l'auto-génération (\texttt{/coordinator/schedules/generation}) pour peupler automatiquement le planning.
    \item Sauvegarder le planning (\texttt{/coordinator/schedules}).
    \item Publier le planning (\texttt{/coordinator/schedules/publication}) pour le rendre visible aux enseignants et étudiants.
\end{enumerate}

\textbf{Suivi et évaluation :}
\begin{itemize}
    \item Consulter les notes (\texttt{/coordinator/grades}) pour suivre l'avancement des évaluations.
    \item Générer et télécharger les documents (convocations, fiches d'évaluation, listes de présence, PV) via la page Documents.
    \item Geler la session (\texttt{freeze}) après la date limite pour bloquer toute modification de notes.
\end{itemize}

\subsection{Enseignant}

\textbf{Déclaration d'indisponibilités :}
\begin{enumerate}
    \item Accéder à la page d'indisponibilités (\texttt{/teacher/unavailability}).
    \item Sélectionner les créneaux d'indisponibilité sur le calendrier interactif.
    \item Sauvegarder. Ces créneaux seront pris en compte par le moteur de conflits (RG6) lors de la planification.
\end{enumerate}

\textbf{Saisie des évaluations :}
\begin{enumerate}
    \item Accéder à la page d'évaluations (\texttt{/teacher/evaluations}).
    \item Sélectionner la soutenance à évaluer.
    \item Saisir le score (sur 20) et un commentaire.
    \item Soumettre l'évaluation. Une fois soumise, la note ne peut plus être modifiée.
\end{enumerate}

\textbf{Consultation du planning :}
\begin{itemize}
    \item La page de planning (\texttt{/teacher/schedule}) affiche les soutenances auxquelles l'enseignant est assigné en tant que membre de jury.
\end{itemize}

\subsection{Étudiant}

\textbf{Gestion du groupe :}
\begin{enumerate}
    \item Accéder à la page du groupe (\texttt{/student/group}).
    \item Consulter les composition des groupes.
    \item Créer un groupe ou rejoindre un groupe existant
\end{enumerate}

\textbf{Consultation des informations :}
\begin{itemize}
    \item Le tableau de bord (\texttt{/student/dashboard}) résume les informations essentielles : soutenance assignée, groupe, documents.
    \item La page de documents (\texttt{/student/documents}) permet de consulter et télécharger les documents requis.
    \item La convocation peut être téléchargée directement depuis l'application.
\end{itemize}
